\documentclass[11pt]{article}
\usepackage{fontspec}
\usepackage[french]{babel}
\usepackage{amsmath}
\usepackage{amssymb}
\usepackage{amsthm}
\usepackage{mathrsfs}
\usepackage{tikz-cd}
\usepackage{enumitem}

\newcommand{\printedpage}[1]{\marginpar{\footnotesize p.~#1}}
\newcommand{\msp}[1]{\marginpar{\footnotesize [#1]}}
\newcommand{\Fltr}{\operatorname{Fl}\operatorname{tr}}
\newcommand{\Ob}{\operatorname{Ob}}
\newcommand{\bDelta}{\boldsymbol{\Delta}}

\begin{document}

% ===================== page 1 =====================
\printedpage{}
% head: (none; title page, no page number printed)

\begin{center}
{\Large\textsf{LES DÉRIVATEURS}}

\bigskip

\textbf{ALEXANDRE GROTHENDIECK}

\bigskip\bigskip

\textbf{Chapitre VII}

\medskip

\textbf{Catégories de chemins (2)}
\end{center}

\newpage

% ===================== page 2 =====================
\printedpage{}
% head: (none; no page number printed)

Ce texte a été déchiffré et transcrit en \LaTeXe\ par M. Künzer. Il a été édité par M. Künzer, J. Malgoire, et G. Maltsiniotis. La transcription est aussi fidèle que possible au manuscrit. Pour les quelques corrections évidentes, ou rares commentaires des éditeurs, ainsi que pour la numérotation originale des pages du manuscrit, les caractères de machine à écrire \texttt{[typewriter]} entre crochets sont utilisés. Un point d'interrogation entre crochets signifie que l'on n'est pas sûr du mot qui précède.

Cette édition est provisoire. Les remarques, commentaires et corrections sont bienvenus. Envoyer un message à :

\textbf{maltsin@math.jussieu.fr}

\hfill G. Maltsiniotis\hspace*{3em}

\newpage

% ===================== page 3 =====================
\printedpage{1}
% head: A. Grothendieck : Dérivateurs, Ch. VII, Édition des Universités de Montpellier II et Paris VII
% foot: Version 22 décembre 2001

\msp{page 1}
Comme le formalisme ‘simplifié’ des chemins que j'ai voulu dégager à la va-vite s'est révélé d'une lourdeur insupportable, de nature à le rendre quasiment inutilisable, je vais essayer de dégager un formalisme souple, collant à l'intuition, au prix seulement d'un effort conceptuel supplémentaire.

\section*{1\quad Intervalles et leur variantes} % [sic] leur variantes -> leurs variantes, plural agreement

On appelle \textit{catégorie intervalle standard} une catégorie définie par un ensemble ordonné, dont l'ensemble des élements est un intervalle d'entiers $[0,n]$ ($n\in\mathbf{N}$, on n'exclut pas le cas $n=0$, i.e. de la catégorie ponctuelle $\{0\}$) : % [sic] élements -> éléments, missing accent
\[
\Ob I=[0,n]\subseteq\mathbf{N},
\]
et satisfaisant l'hypothèse suivante :

\begin{enumerate}
\item[a)] pour toute paire de deux sommets consécutifs $i-1$, $i$ de $I$ ($1\leq i\leq n$), il y a une flèche qui les relie, donc soit $i\longrightarrow i+1$, soit $i\longleftarrow i+1$ \texttt{[plutôt} $i-1\longrightarrow i$\texttt{,} $i-1\longleftarrow i$\texttt{]} (cette flèche est notée $\underline{u}_i$,
\[
(i-1)\xrightarrow{\underline{u}_i}i\qquad\quad\text{ou}\qquad\quad(i-1)\xleftarrow{\underline{u}_i}i\ \ ).
\]
\item[b)] La catégorie $I$ est engendrée par l'ensemble de ses flèches $\underline{u}_i$ (appelées \textit{flèches de transition}).
\end{enumerate}

La flèche de transition est dite \textit{directe} si elle va de $i-1$ à $i$, \textit{rétrograde} dans le sens inverse. Si on affecte l'intervalle $[i-1,i]$ du signe $+$ ou $-$ suivant que $\underline{u}_i$ est directe ou rétrograde, on voit donc que les
\msp{page 2}
intervalles standard de longueur donnée $n$ correspondent biunivoquement aux systèmes de signes $+$ ou $-$ sur les $n$ intervalles de transition $[i-1,i]$ ($1\leq i\leq n$), il y en a donc $2^n$. Le sommet $0$ est l'\textit{origine} ou la \textit{source} de l'intervalle, le sommet $n$ l'\textit{extrémité} ou le \textit{but} de l'intervalle.

Un intervalle standard est \textit{orienté,} du fait qu'on à la relation d'ordre \textit{total} naturelle sur l'ensemble de ses sommets, savoir celle des entiers dans $[0,n]$. Quand on parlera de la relation d'ordre sur $I$, qu'un sommet $i$ en \textit{‘précède’} ou \textit{‘suit’} un autre $j$ ($i\leq j$ ou $i\geq j$ respectivement), c'est de cette relation d'ordre qu'on voudra parler. Celle associée à la structure de catégorie ne sera jamais regardée comme une relation d'ordre. Il y a également une relation d'ordre total naturelle sur l'ensemble des flèches de transition $\underline{u}_i$ ($1\leq i\leq n$), donnée par la relation d'ordre des entiers $i$, $1\leq i\leq n$, qui les indexent. On dira qu'une flèche $i$ \texttt{[plutôt} $\underline{u}_i$\texttt{]} de transition \textit{précède} ou \textit{suit} une autre \texttt{[}$\underline{u}_j$\texttt{]}, ou qu'elle est \textit{successeur immédiat} ou \textit{prédécesseur immédiat} d'une autre (si $i=j+1$ ou $i=j-1$). % [sic] qu'on à -> qu'on a, verb; [sic] un sommet i en ‘précède’ -> un sommet i ‘précède’, superfluous "en"

\newpage

% ===================== page 4 =====================
\printedpage{2}
% head: A. Grothendieck : Dérivateurs, Ch. VII, Édition des Universités de Montpellier II et Paris VII
% foot: Version 22 décembre 2001

Si on a un ensemble d'entiers
\[
E\subseteq\mathbf{Z},
\]
on considérera dans $E$ la relation d'équivalence engendrée par la relation entre deux éléments $(i,j)$ de $E$ : $j$ est successeur immédiat de $i$ \textit{dans} $\mathbf{Z}$ (i.e. $j=i+1$). Deux éléments distinct $i$, $j$ de $E$ sont congrus pour cette relation si et seulement si l'intervalle d'entiers qu'ils déterminent est $\subseteq E$, donc % [sic] distinct -> distincts, plural agreement
\[
\text{si } i<j,\ i \text{ et } j \text{ sont \textit{‘connectés’} dans } E \text{ si et seulement si } [i,j]\subseteq E.
\]

\msp{page 3}
Quand on parlera des ‘composantes connexes’ de $E$, il s'agira toujours des classes suivant cette relation d'équivalence.

Nous identifierons l'ensemble des flèches de transition $\Fltr(I)$ de $I$ à l'intervalle d'entiers $[1,n]$, donc à une partie de $\mathbf{Z}$. Si on a une partie $E\subseteq\Fltr(I)$, on y a donc la notion de flèches de transition ‘connectées’ dans $E$, et la notion de connexité pour $E$. Si $i<j$, $\underline{u}_i$ et $\underline{u}_j$ sont connectées si et seulement si toutes les flèches de transition intermédiaires entre $\underline{u}_i$ et $\underline{u}_j$ sont dans $E$. Et $E$ est connexe si et seulement s'il est soit vide, soit formée d'une suite de flèches de transition dont chacune est successeur immédiat de celle qui la précède, % [sic] soit formée -> soit formé, agreement with "il"
\[
E=\{\underline{u}_i,\underline{u}_{i+1}\ldots,\underline{u}_{i+l-1}\}, % [sic] \underline{u}_{i+1}\ldots -> \underline{u}_{i+1},\ldots, missing comma
\]
où $l$ est le cardinal de $E$. Une telle $E$ est appelée une \textit{chaîne} de flèches de transition, et $l$ est sa longueur. Une chaîne est dite \textit{composable} si le composé $\underline{u}_{i+l-1}\circ\cdots\circ\underline{u}_{i+1}\circ\underline{u}_i$ \texttt{[ou} $\underline{u}_i\circ\underline{u}_{i+1}\circ\ldots\circ\underline{u}_{i+l-1}$\texttt{]} a un sens, i.e. si elles sont toutes de même sens, soit directe, soit rétrograde.

On désigne par $\Fltr^+(I)$, $\Fltr^-(I)$ l'ensemble des flèches de transition directes, ou rétrogrades, respectivement. Leurs composantes connexes ne sont autres que les chaînes composables maximales. Ainsi, les chaînes composables maximales forment une partition de l'ensemble des flèches de transition, en intervalles d'entiers mutuellement disjoints. Les chaînes forment elles-mêmes un ensemble totalement ordonné
\[
C_1,\ C_2,\ \ldots\ C_r\qquad\quad(1\leq r\leq n).
\]
L'entier $r$ s'appelle la \textit{‘longueur utile’} de l'intervalle
\msp{page 4}
$I$, c'est donc le cardinal de l'ensemble de ses chaînes composables maximales.

Deux chaînes composables maximales \texttt{[consécutives]} sont toujours de signe opposé, si l'une est directe l'autre est rétrograde, et inversement. La première chaîne composable maximale $C_1$ (disons simplement \textit{‘chaîne structurale’}) a le même signe que $\underline{u}_1$; \texttt{[on]} dit que la \textit{direction initiale} de $I$ est directe, resp. rétrograde, s'il en est ainsi de $\underline{u}_1$, ou ce qui revient au même, de $C_1$. De même on définit la \textit{direction finale} de $I$ comme celle de $\underline{u}_n$, ou encore celle de $C_r$. Direction initiale et finale sont égales si la longueur utile $r$ est
\printedpage{3}
% head: A. Grothendieck : Dérivateurs, Ch. VII, Édition des Universités de Montpellier II et Paris VII
% foot: Version 22 décembre 2001
impaire, inégales si elle est paire. (\textbf{NB} Ce n'est que pour $n=0$, i.e. quand $I=\{0\}$ est la catégorie ponctuelle, que direction initiale et direction finale ne sont pas définies.)

Les chaînes composables correspondent aux sous-catégories pleines de $I$ qui sont isomorphes à une catégorie $\bDelta_i$ (où $i$ est la longueur de la chaîne). Notons que $\bDelta_i$ n'a pas d'autre automorphisme que l'identité, donc les chaînes maximales de longueur donnée $i\geq 1$ correspondent exactement aux foncteurs
\[
\bDelta_i\lhook\joinrel\longrightarrow\ I
\]
qui sont injectifs sur les sommets. Un tel foncteur est pleinement fidèle, et un monomorphisme. Ce foncteur définit une application strictement croissante ou strictement décroissante sur
\msp{page 5}
les sommets, suivant que la chaîne est directe ou rétrograde.

Notons que tout simplexe type $\bDelta_i$, $i\geq 0$, peut être vu comme un intervalle, de longueur $i$, de longueur utile $1$ (si $i\geq 1$, $0$ si $i=0$). Son origine est $0$, son extrémité est $i$. Mais appelons une catégorie quelconque une \textit{catégorie intervalle} si elle est isomorphe à un intervalle standard. Notons ici

\textbf{Proposition 1.} \textit{Soit $I$ un intervalle standard de longueur $n\geq 1$.}

\begin{enumerate}
\item[\textit{a)}] \textit{Soit $i\in I$. Pour que $i\in\{0,n\}$, i.e. pour que $i$ soit un des deux} ‘bouts’ \textit{(ou} ‘extrémités’\textit{) (source ou but) de $I$, il faut et il suffit que l'une des deux catégories $I/i$ ou $i\backslash I$ soit la catégorie ponctuelle et que l'autre soit isomorphe à un} $\bDelta_j$ (\footnote{où $j$ est la longueur de l'unique chaîne structurale qui ‘touche’ $i$ (i.e. telle que $i$ soit dans la sous-catégorie pleine définie par ladite chaîne).}). \textit{(Caractérisation intrinsèque des bouts d'un intervalle.) Pour que $i$ soit la source, il faut et il suffit que $i\backslash I$} \texttt{[phrase incomplète]}
\item[\textit{b)}] \textit{Tout foncteur $I\longrightarrow I$ qui fixe les deux extrémités de $I$ est l'identité.}
\end{enumerate}

\textbf{Corollaire.} \textit{Le groupe des automorphismes de $I$ a au plus deux éléments. S'il en a deux, l'automorphisme non trivial échange les deux extrémités. (On dit alors que l'intervalle est} symétrique.\textit{)}

\msp{page 6}
NB Si $i$ n'est pas un bout, et si néanmoins l'une des deux catégories $i\backslash I$, $I/i$ est ponctuelle, disons $I/i$, alors au voisinage immédiat de $i$, $I$ a la forme
\[
\begin{tikzcd}[row sep=0pt]
\bullet & \bullet \arrow[l, "\underline{u}_i"'] \arrow[r, "\underline{u}_{i+1}"] & \bullet \\
i-1 & i & i+1
\end{tikzcd}
\]
($\underline{u}_i$ rétrograde, $\underline{u}_{i+1}$ directe - si c'était $i\backslash I$ la ponctuelle, alors $\underline{u}_i$ serait directe, $\underline{u}_{i+1}$ rétrograde - en tous les cas ces deux flèches sont de sens opposé, i.e. non composables). On trouve donc que $i\backslash I$ est une catégorie de la forme

\end{document}
